\begin{lemma}\label{lem:divisibility}
  Let \(n \in \mathbb{Z}\). If \(7 \mid n^2\), then \(7 \mid n\).
\end{lemma}

\begin{proof}
  By the division algorithm, applied with positive divisor \(7\), there exist
  \(q \in \mathbb{Z}\) and \(r \in \{0,1,2,3,4,5,6\}\) such that
  \[
    n = 7q + r.
  \]
  We first justify explicitly that \(n^2\) and \(r^2\) have the same remainder
  modulo \(7\). Indeed,
  \[
    n^2 - r^2 = (7q+r)^2 - r^2
    = 49q^2 + 14qr
    = 7(7q^2+2qr).
  \]
  Since \(7q^2+2qr \in \mathbb{Z}\), this shows that
  \[
    7 \mid (n^2-r^2).
  \]

  Now assume \(7 \mid n^2\). Since both \(n^2\) and \(n^2-r^2\) are divisible
  by \(7\), their difference is also divisible by \(7\). Hence
  \[
    r^2 = n^2 - (n^2-r^2)
  \]
  is divisible by \(7\).

  We now check the possible values of \(r\):
  \[
  \begin{array}{c|ccccccc}
    r & 0 & 1 & 2 & 3 & 4 & 5 & 6 \\
    \hline
    r^2 & 0 & 1 & 4 & 9 & 16 & 25 & 36 \\
    r^2 \bmod 7 & 0 & 1 & 4 & 2 & 2 & 4 & 1
  \end{array}
  \]
  Thus, among \(r \in \{0,1,2,3,4,5,6\}\), the condition \(7 \mid r^2\) holds
  only when \(r=0\). Therefore \(r=0\), and so
  \[
    n = 7q.
  \]
  Hence \(7 \mid n\).
\end{proof}

\begin{theorem}\label{thm:sqrt7}
  \(\sqrt{7}\) is irrational.
\end{theorem}

\begin{proof}
  Suppose, for contradiction, that \(\sqrt{7}\) is rational. Then there exist
  integers \(m,n\in\mathbb{Z}\) with \(n\neq 0\) such that
  \[
    \sqrt{7} = \frac{m}{n}.
  \]
  Since \(n\neq 0\), either \(n>0\) or \(n<0\). If \(n<0\), replace \(m\) and
  \(n\) by \(-m\) and \(-n\). This does not change the fraction because
  \[
    \frac{-m}{-n} = \frac{m}{n},
  \]
  and the new denominator \(-n\) is positive. Therefore we may assume
  \[
    n>0.
  \]

  Let
  \[
    g = \gcd(m,n).
  \]
  Since \(n\neq 0\), the integers \(m\) and \(n\) are not both zero, so
  \(g\) is a positive integer. In particular,
  \[
    g \neq 0.
  \]
  Because \(g\mid m\) and \(g\mid n\), there exist integers \(a,b\in\mathbb{Z}\)
  such that
  \[
    m=ga, \qquad n=gb.
  \]
  Since \(n>0\) and \(g>0\), we have \(b>0\), so in particular \(b\neq 0\).
  Therefore
  \[
    \sqrt{7}
    = \frac{m}{n}
    = \frac{ga}{gb}
    = \frac{a}{b},
  \]
  where the cancellation is valid because \(g\neq 0\).

  We claim that \(\gcd(a,b)=1\). Let \(d\) be any positive integer dividing both
  \(a\) and \(b\). Then \(a=du\) and \(b=dv\) for some integers \(u,v\). Hence
  \[
    m=ga=gdu, \qquad n=gb=gdv,
  \]
  so \(gd\) is a positive common divisor of \(m\) and \(n\). Since
  \(g=\gcd(m,n)\) is the greatest positive common divisor of \(m\) and \(n\),
  we must have
  \[
    gd \leq g.
  \]
  Because \(g>0\), dividing by \(g\) gives \(d\leq 1\). Since \(d\) is a positive
  integer, \(d=1\). Thus the only positive common divisor of \(a\) and \(b\) is
  \(1\), so
  \[
    \gcd(a,b)=1.
  \]

  Now square both sides of \(\sqrt{7} = a/b\).
  Since equality is preserved under squaring,
  \[
    (\sqrt{7})^2 = \left(\frac{a}{b}\right)^2.
  \]
  By the definition of \(\sqrt{7}\), we have \((\sqrt{7})^2 = 7\).
  Also, since \(b\neq 0\),
  \[
    \left(\frac{a}{b}\right)^2 = \frac{a^2}{b^2}.
  \]
  Therefore \(7 = a^2/b^2\). Multiplying both sides by \(b^2\) gives
  \[
    a^2 = 7b^2.
  \]

  Hence \(7 \mid a^2\). Since \(a\in\mathbb{Z}\), Lemma~\ref{lem:divisibility}
  applies, and we get \(7 \mid a\).
  Thus \(a=7k\) for some \(k\in\mathbb{Z}\). Substituting into \(a^2=7b^2\):
  \[
    (7k)^2 = 7b^2 \implies 49k^2 = 7b^2 \implies 7k^2 = b^2.
  \]
  Thus \(7 \mid b^2\). Since \(b\in\mathbb{Z}\), Lemma~\ref{lem:divisibility}
  applies again, and therefore \(7 \mid b\).

  We have shown that \(7\mid a\) and \(7\mid b\). Hence \(7\) is a positive
  common divisor of \(a\) and \(b\), with \(7>1\), contradicting
  \(\gcd(a,b)=1\).

  Therefore our assumption that \(\sqrt{7}\) is rational is false. Hence
  \(\sqrt{7}\) is irrational.
\end{proof}
